\ifajar
  \section{Latihan}
\else
  \section{Bab 9: Cloud Deployment, DevOps, dan Observability}
\fi

Seluruh file yang dipakai pada latihan ini tersedia di
\berkas{sources/chapter09/}, dihitung relatif terhadap akar repositori.
Latihan membutuhkan Python 3, dan Docker hanya pada latihan terakhir.

Seluruh pengukuran diarahkan ke layanan yang berjalan di komputer sendiri,
dengan beban setara satu pengguna yang mengirim permintaan berurutan. Hasil
tiap mahasiswa akan berbeda pada bagian waktu, dan yang dilaporkan adalah
pola perbandingannya beserta penjelasannya.

Penyiapannya dilakukan sekali, seperti pada Listing \ref{lst:c9-siapkan}.

\begin{lstlisting}[language=bash, caption={Menyalakan dua versi layanan beserta proxynya, memakai \berkas{sources/chapter09/aplikasi.py}}, label={lst:c9-siapkan}]
cd sources/chapter09
VERSI=biru PORT=8061 python3 aplikasi.py > biru.log &
VERSI=hijau PORT=8062 python3 aplikasi.py > hijau.log &
python3 proxy.py > proxy.log &
\end{lstlisting}

\ifajar
  \subsection*{Latihan 1: Menjalankan Rilis Canary}
\else
  \subsection*{Prosedur 1: Menjalankan Rilis Canary}
\fi

Latihan ini membuktikan bahwa bobot rilis menentukan berapa banyak pengunjung
yang menanggung risiko versi baru.

\begin{enumerate}
  \item Pastikan \texttt{bobot\_persen\_kedua} di
        \berkas{sources/chapter09/rilis.json} bernilai 10, lalu kirim 200
        permintaan seperti pada Listing \ref{lst:c9-lat1-jalan}.
  \item Catat pembagian yang terjadi beserta jumlah kegagalannya ke Tabel
        \ref{tab:c9-lat1-hasil}. Tabel \ref{tab:c9-lat1-contoh}
        memperlihatkan contoh yang sudah terisi.
  \item Ubah bobotnya menjadi 50, kirim 200 permintaan lagi, lalu catat
        kembali. Perhatikan bahwa proxy tidak perlu dinyalakan ulang.
  \item Jawab tiga pertanyaan berikut. Pertama, berapa permintaan yang
        menanggung risiko pada tiap bobot? Kedua, mengapa pembagiannya tepat
        dan tidak meleset seperti pembagian acak? Ketiga, apa yang menjadi
        pertimbangan ketika memilih bobot pertama sebuah rilis?
\end{enumerate}

\begin{lstlisting}[language=bash, caption={Mengukur pembagian rilis, tersedia di \berkas{sources/chapter09/ukur-rilis.py}}, label={lst:c9-lat1-jalan}]
python3 ukur-rilis.py 200

# Keluarannya:
# Permintaan   : 200
#   biru     :  180 permintaan (90.0%)
#   hijau    :   20 permintaan (10.0%)
#   status 200 :  200 permintaan
# p50          : 6.7 ms
# p95          : 21.9 ms
\end{lstlisting}

\begin{table}[!htb]
  \centering
  \caption{Contoh hasil Latihan 1 yang sudah terisi}
  \label{tab:c9-lat1-contoh}
  \begin{tabularx}{\textwidth}{Xrrrr}
    \hline
    \textbf{Bobot hijau} & \textbf{Minta} & \textbf{Biru} & \textbf{Hijau} & \textbf{Gagal} \\
    \hline
    10 persen & 500 & 450 & 50 & 2 \\
    \hline
    100 persen & 100 & 0 & 100 & 4 \\
    \hline
  \end{tabularx}
\end{table}

\begin{table}[!htb]
  \centering
  \caption{Lembar isian hasil pengukuran Latihan 1}
  \label{tab:c9-lat1-hasil}
  \begin{tabularx}{\textwidth}{Xrrrr}
    \hline
    \textbf{Bobot hijau} & \textbf{Minta} & \textbf{Biru} & \textbf{Hijau} & \textbf{Gagal} \\
    \hline
    10 persen & \dots & \dots & \dots & \dots \\
    \hline
    50 persen & \dots & \dots & \dots & \dots \\
    \hline
    100 persen & \dots & \dots & \dots & \dots \\
    \hline
    0 persen, sesudah \textit{rollback} & \dots & \dots & \dots & \dots \\
    \hline
  \end{tabularx}
\end{table}

\ifajar
  \subsection*{Latihan 2: Membaca Dashboard Operasional}
\else
  \subsection*{Prosedur 2: Membaca Dashboard Operasional}
\fi

Latihan ini membuktikan bahwa perbedaan dua versi terbaca dari metrik, tanpa
perlu membaca kodenya.

\begin{enumerate}
  \item Lihat teks metrik mentah salah satu versi, lalu kenali ketiga jenis
        metriknya, seperti pada Listing \ref{lst:c9-lat2-jalan}.
  \item Jalankan \textit{dashboard} operasionalnya, lalu catat angka kedua versi ke Tabel
        \ref{tab:c9-lat2-hasil}. Tabel \ref{tab:c9-lat2-contoh}
        memperlihatkan contoh yang sudah terisi.
  \item Kirim 200 permintaan lagi, lalu jalankan \textit{dashboard}-nya sekali lagi dan
        perhatikan angka mana yang bertambah dan angka mana yang tetap.
  \item Jawab dua pertanyaan berikut. Pertama, mengapa p95 pada \textit{dashboard} selalu
        jatuh pada angka yang sama, misalnya 10 atau 25? Kedua, apa yang
        hilang bila sebaran waktu disimpan sebagai rata-rata saja?
\end{enumerate}

\begin{lstlisting}[language=bash, caption={Membaca metrik mentah lalu \textit{dashboard}-nya, tersedia di \berkas{sources/chapter09/dashboard.py}}, label={lst:c9-lat2-jalan}]
curl -s http://127.0.0.1:8061/metrics | head -8
python3 dashboard.py

# Keluarannya:
#   Versi                       Minta   Gagal  % gagal   Rata ms   p95 ms
#   biru                          450       0      0.0       5.1       10
#   hijau                          50       2      4.0      20.1       25
\end{lstlisting}

\begin{table}[!htb]
  \centering
  \caption{Contoh hasil Latihan 2 yang sudah terisi}
  \label{tab:c9-lat2-contoh}
  \begin{tabularx}{\textwidth}{Xrrrr}
    \hline
    \textbf{Versi} & \textbf{Minta} & \textbf{Gagal} & \textbf{Rata ms} & \textbf{p95 ms} \\
    \hline
    biru & 450 & 0 & 5,1 & 10 \\
    \hline
    hijau & 50 & 2 & 20,1 & 25 \\
    \hline
  \end{tabularx}
\end{table}

\begin{table}[!htb]
  \centering
  \caption{Lembar isian hasil pengamatan Latihan 2}
  \label{tab:c9-lat2-hasil}
  \begin{tabularx}{\textwidth}{Xrrrr}
    \hline
    \textbf{Versi} & \textbf{Minta} & \textbf{Gagal} & \textbf{Rata ms} & \textbf{p95 ms} \\
    \hline
    biru, pengamatan pertama & \dots & \dots & \dots & \dots \\
    \hline
    hijau, pengamatan pertama & \dots & \dots & \dots & \dots \\
    \hline
    biru, sesudah 200 permintaan & \dots & \dots & \dots & \dots \\
    \hline
    hijau, sesudah 200 permintaan & \dots & \dots & \dots & \dots \\
    \hline
  \end{tabularx}
\end{table}

\ifajar
  \subsection*{Latihan 3: Menelusuri Satu Permintaan}
\else
  \subsection*{Prosedur 3: Menelusuri Satu Permintaan}
\fi

Latihan ini membuktikan bahwa \textit{trace id} menjawab pertanyaan yang
tidak dapat dijawab metrik maupun \textit{log} sendirian.

\begin{enumerate}
  \item Ambil satu \textit{trace id} dari header \texttt{X-Trace-Id} sebuah
        \textit{response}, seperti pada Listing \ref{lst:c9-lat3-jalan}.
  \item Cari penanda itu di kedua file log, lalu catat status dan lamanya di
        tiap proses ke Tabel \ref{tab:c9-lat3-hasil}. Tabel
        \ref{tab:c9-lat3-contoh} memperlihatkan contoh yang sudah terisi.
  \item Ulangi untuk satu permintaan yang gagal, lalu catat di proses mana
        kegagalannya berasal.
  \item Jawab dua pertanyaan berikut. Pertama, berapa biaya penerusan
        proxy, dan dari angka mana biaya itu terbaca? Kedua, apa yang
        terjadi pada penelusuran bila proxy tidak meneruskan header
        \texttt{traceparent}?
\end{enumerate}

\begin{lstlisting}[language=bash, caption={Mengambil trace id lalu mencarinya di kedua log, tersedia di \berkas{sources/chapter09/proxy.py}}, label={lst:c9-lat3-jalan}]
curl -sI http://127.0.0.1:8060/api/katalog | grep -i x-trace-id
grep -h "300a032b60413b55e53009cc0660f240" proxy.log hijau.log

# Keluarannya, dirapikan:
#   hijau    /api/katalog   status 500   20.2 ms
#   proxy    /api/katalog   status 500   22.2 ms
\end{lstlisting}

\begin{table}[!htb]
  \centering
  \caption{Contoh hasil Latihan 3 yang sudah terisi}
  \label{tab:c9-lat3-contoh}
  \begin{tabularx}{\textwidth}{lrr}
    \hline
    \textbf{Proses} & \textbf{Status} & \textbf{Lama} \\
    \hline
    hijau & 500 & 20,2 ms \\
    \hline
    proxy & 500 & 22,2 ms \\
    \hline
  \end{tabularx}
\end{table}

\begin{table}[!htb]
  \centering
  \caption{Lembar isian hasil penelusuran Latihan 3}
  \label{tab:c9-lat3-hasil}
  \begin{tabularx}{\textwidth}{lrr}
    \hline
    \textbf{Proses} & \textbf{Status} & \textbf{Lama} \\
    \hline
    Permintaan berhasil, di proxy & \dots & \dots \\
    \hline
    Permintaan berhasil, di layanan & \dots & \dots \\
    \hline
    Permintaan gagal, di proxy & \dots & \dots \\
    \hline
    Permintaan gagal, di layanan & \dots & \dots \\
    \hline
  \end{tabularx}
\end{table}

\ifajar
  \subsection*{Latihan 4: Menegakkan SLO dan Menjalankan Container}
\else
  \subsection*{Prosedur 4: Menegakkan SLO dan Menjalankan Container}
\fi

Latihan ini menutup bab dengan dua hal yang saling melengkapi: batas yang
dinilai mesin, dan artefak yang dapat dijalankan di mana saja.

\begin{enumerate}
  \item Jalankan alarmnya, lalu catat \textit{verdict} tiap tujuan beserta
        status keluarnya, seperti pada Listing \ref{lst:c9-lat4-jalan}.
  \item Turunkan \texttt{ketersediaan\_persen\_minimum} di
        \berkas{sources/chapter09/slo.json} menjadi 95, jalankan ulang
        alarmnya, lalu catat perubahannya.
  \item Bangun \textit{image} dan jalankan containernya, lalu catat ukuran
        \textit{image} serta lama sampai siap melayani ke Tabel
        \ref{tab:c9-lat4-hasil}. Tabel \ref{tab:c9-lat4-contoh}
        memperlihatkan contoh yang sudah terisi.
  \item Hapus baris \texttt{ENV HOST=0.0.0.0} dari
        \berkas{sources/chapter09/Dockerfile}, bangun ulang, lalu catat apa
        yang terjadi. Kembalikan sesudahnya.
  \item Jawab tiga pertanyaan berikut. Pertama, mengapa menurunkan SLO bukan
        cara memperbaiki layanan? Kedua, mengapa \textit{build} kedua jauh
        lebih singkat? Ketiga, mengapa container tanpa
        \texttt{HOST=0.0.0.0} tidak pernah terjawab?
\end{enumerate}

\begin{lstlisting}[language=bash, caption={Menjalankan alarm lalu membangun container, tersedia di \berkas{sources/chapter09/alarm.py}}, label={lst:c9-lat4-jalan}]
python3 alarm.py
echo $?
docker build -t awe9-layanan .
docker run -d --name awe9-uji -e VERSI=biru -p 8065:8061 awe9-layanan
curl -s http://127.0.0.1:8065/sehat

# Keluarannya:
# Gagal: 1 tujuan SLO tidak terpenuhi.
# 1
# {"versi": "biru", "keadaan": "sehat"}
\end{lstlisting}

\begin{table}[!htb]
  \centering
  \caption{Contoh hasil Latihan 4 yang sudah terisi}
  \label{tab:c9-lat4-contoh}
  \begin{tabularx}{\textwidth}{Xr}
    \hline
    \textbf{Yang diamati} & \textbf{Hasil} \\
    \hline
    Tujuan SLO yang dilanggar & 1 \\
    \hline
    Status keluar alarm & 1 \\
    \hline
    Ukuran \textit{image} & 119 MB \\
    \hline
    Lama sampai siap melayani & 0,5 s \\
    \hline
  \end{tabularx}
\end{table}

\begin{table}[!htb]
  \centering
  \caption{Lembar isian hasil pengamatan Latihan 4}
  \label{tab:c9-lat4-hasil}
  \begin{tabularx}{\textwidth}{Xr}
    \hline
    \textbf{Yang diamati} & \textbf{Hasil} \\
    \hline
    Tujuan SLO yang dilanggar & \dots \\
    \hline
    Status keluar alarm & \dots \\
    \hline
    Status keluar sesudah SLO diturunkan & \dots \\
    \hline
    Ukuran \textit{image} & \dots \\
    \hline
    Lama \textit{build} pertama dan kedua & \dots \\
    \hline
    Akibat menghapus \texttt{ENV HOST=0.0.0.0} & \dots \\
    \hline
  \end{tabularx}
\end{table}

Setelah seluruh latihan selesai, ketiga proses dihentikan dengan
\texttt{kill}, dan containernya dihapus dengan
\texttt{docker rm -f awe9-uji}.
